% !TEX root = ../../hw07.tex
% Source: author's supplied text in chat, 2026-11-06, Q2 (i)--(ii).
% Preserve the two squeeze arguments and the one-sided quotient calculations.
\begin{proof}[Solution]
  Part (i). Recall that continuity at $c$ is defined by
  \[
    \lim_{x\to c}f(x)=f(c).
  \]
  We want to prove that $\lim_{x\to0}f(x)=0=f(0)$. Analyzing both
  left and right sides, where $x\to0^-$ and $x\to0^+$ respectively,
  we have
  \[
    0\le\frac{x}{1+e^{1/x}}\le x
    \qquad(x>0).
  \]
  By the squeeze theorem,
  \[
    \lim_{x\to0^+}\frac{x}{1+e^{1/x}}=0.
  \]
  For $x\to0^-$ the inequality is reversed,
  \[
    x\le\frac{x}{1+e^{1/x}}\le0,
  \]
  and similarly, by the squeeze theorem,
  \[
    \lim_{x\to0^-}\frac{x}{1+e^{1/x}}=0.
  \]
  Therefore, since
  \[
    \lim_{x\to0^-}\frac{x}{1+e^{1/x}}
      =\lim_{x\to0^+}\frac{x}{1+e^{1/x}}=0,
  \]
  $f$ is continuous at $x=0$.

  For part (ii), recall that to find $f_-'(0)$ and $f_+'(0)$ we
  use the definitions of the left- and right-side derivatives,
  \[
    \begin{aligned}
      f_+'(c)&=\lim_{h\to0^+}\frac{f(c+h)-f(c)}{h},\\
      f_-'(c)&=\lim_{h\to0^-}\frac{f(c+h)-f(c)}{h}.
    \end{aligned}
  \]
  \autoref{fig:hw07:one-sided-slopes} shows why
  continuity can hold even though these two slopes differ.
  \begin{marginfigure}
    \centering
    \input{\analysisroot/figures/interval-styles.tex}
\begin{tikzpicture}[interval diagram,x=1.65cm,y=1.65cm]
  \draw[->] (-1.1,0)--(1.22,0) node[right=4pt] {$x$};
  \draw[->] (0,-.93)--(0,.52) node[above=4pt] {$y$};
  \node[below right=5pt] at (0,0) {$0$};
  \draw[gray,dashed] (-.86,-.86)--(-.08,-.08);
  \node[below right=6pt] at (-.9,-.68) {$y=x$};
  \draw[gray,dashed] (.12,0)--(1.05,0);
  \node[below=8pt] at (.8,0) {$y=0$};
  \draw[<-,domain=-1.03:-.025,samples=80]
    plot (\x,{\x/(1+exp(1/\x))});
  \draw[->,domain=.025:1.05,samples=80]
    plot (\x,{\x*exp(-1/\x)/(1+exp(-1/\x))});
  \fill (0,0) circle[radius=1.6pt];
  \node[above=6pt] at (.78,{.78/(1+exp(1/.78))}) {$f$};
  \draw (-1.02,-1.17)--(-.72,-1.17);
  \node[anchor=west] at (-.6,-1.17) {function};
  \draw[gray,dashed] (-1.02,-1.43)--(-.72,-1.43);
  \node[anchor=west] at (-.6,-1.43) {one-sided tangents};
\end{tikzpicture}

    \caption[Continuity with unequal one-sided slopes.]{%
      The graph meets the origin continuously. Its left tangent is
      $y=x$, while its right tangent is $y=0$. The two limiting
      slopes are $1$ and $0$.}
    \label{fig:hw07:one-sided-slopes}
  \end{marginfigure}
  Starting with the left limit,
  \[
    \begin{aligned}
      f_-'(0)
        &=\lim_{h\to0^-}\frac{f(h)-f(0)}{h}\\
        &=\lim_{h\to0^-}
          \frac{\frac{h}{1+e^{1/h}}-f(0)}{h}\\
        &=\lim_{h\to0^-}\frac{\frac{h}{1+e^{1/h}}}{h}\\
        &=\lim_{h\to0^-}\frac1{1+e^{1/h}}.
    \end{aligned}
  \]
  Since $h\to0^-$ implies $e^{1/h}\to0$, we have $f_-'(0)=1$.

  Similarly, $f_+'(0)=0$, since
  \[
    h\to0^+
      \quad\Longrightarrow\quad \frac1h\to+\infty
      \quad\Longrightarrow\quad e^{1/h}\to+\infty
      \quad\Longrightarrow\quad \frac1{1+e^{1/h}}\to0.
  \]
  Therefore, since $f_-'(0)\ne f_+'(0)$, $f'(0)$ is undefined.
\end{proof}
